\documentclass[aps,prb,onecolumn,nofootinbib,superscriptaddress]{revtex4-2}
\usepackage{amsmath,amssymb,amsthm,mathtools,bm}
\usepackage{booktabs}
\usepackage[colorlinks=true,linkcolor=blue,citecolor=blue,urlcolor=blue]{hyperref}
\usepackage{microtype}

\newtheorem{definition}{Definition}
\newtheorem{proposition}{Proposition}
\newtheorem{remark}{Remark}

\newcommand{\cH}{\mathcal H}
\newcommand{\cW}{\mathcal W}
\newcommand{\cR}{\mathcal R}
\newcommand{\cY}{\mathcal Y}
\newcommand{\cZ}{\mathcal Z}
\newcommand{\Id}{\mathbf{1}}
\newcommand{\vac}{\mathrm{vac}}
\newcommand{\REF}{\mathrm{REF}}
\newcommand{\full}{\mathrm{full}}
\newcommand{\Ran}{\operatorname{Ran}}
\newcommand{\Ker}{\operatorname{Ker}}
\newcommand{\Span}{\operatorname{span}}
\newcommand{\Tr}{\operatorname{Tr}}
\newcommand{\diag}{\operatorname{diag}}
\newcommand{\bra}[1]{\langle #1\rvert}
\newcommand{\ket}[1]{\lvert #1\rangle}
\newcommand{\braket}[2]{\langle #1\vert #2\rangle}
\newcommand{\ahat}{\hat a}
\newcommand{\chat}{\hat c}
\newcommand{\eps}{\epsilon}

\begin{document}

\title{Occupied Frames, Choi Covariances, and Effective Transfer Matrices for Gain--Loss Strings}
\author{Abid H. Bhuiyan}
\date{\today}

\begin{abstract}
This note gives a frame-level derivation of effective particle and hole
transfer matrices for Choi states generated by ordered strings of
number-conserving Gaussian operators, gains, and losses.  The construction
begins with the one-particle basis and the covariance-space convention for
fermion modes, then rewrites fixed-charge Gaussian states as Slater occupied
frames.  Gain and loss operators act by adding or removing one direction from
the occupied frame.  In doubled Choi space the same operation propagates the
occupied projector of the Choi state itself.  Effective particle and hole
transfer matrices are then graph coordinates of the final occupied or empty
Choi subspace, with exact occupation caps appearing as zeros and their
regulated complements as reciprocal poles.
\end{abstract}

\maketitle
\setcounter{tocdepth}{2}
\tableofcontents

\section{One-Particle Basis and Covariance-Space Modes}

Let \(\cH\simeq \mathbb C^N\) be the one-particle Hilbert space with fixed
orthonormal basis
\begin{equation}
  \{e_i\}_{i=1}^N,\qquad e_i^\dagger e_j=\delta_{ij}.
\end{equation}
The fermion operators \(\chat_i,\chat_i^\dagger\) obey
\begin{equation}
  \{\chat_i,\chat_j^\dagger\}=\delta_{ij},\qquad
  \{\chat_i,\chat_j\}=0,\qquad
  \{\chat_i^\dagger,\chat_j^\dagger\}=0 .
\end{equation}
Throughout this note a one-particle column vector \(x\in\cH\) labels the
covariance-space creation and annihilation modes
\begin{equation}
  \ahat^\dagger(x)=\sum_i x_i^*\chat_i^\dagger,\qquad
  \ahat(x)=\sum_i x_i\chat_i .
  \label{eq:modes}
\end{equation}
The complex conjugation in \(\ahat^\dagger(x)\) is part of the convention.
It gives
\begin{equation}
  \{\ahat(\ell),\ahat^\dagger(g)\}=g^\dagger \ell .
  \label{eq:mode-car}
\end{equation}
Thus the same column \(x\) labels the one-particle state
\(\ahat^\dagger(x)\ket{\vac}\) and its rank-one one-body density matrix
\(xx^\dagger\).

For a nonzero vector \(x\), define the orthogonal projector
\begin{equation}
  P[x]=x(x^\dagger x)^{-1}x^\dagger,\qquad Q[x]=\Id-P[x].
\end{equation}
For a full-rank frame \(X=(x_1,\ldots,x_r)\), define
\begin{equation}
  P[X]=X(X^\dagger X)^{-1}X^\dagger,\qquad Q[X]=\Id-P[X].
  \label{eq:frame-projector}
\end{equation}
This formula is independent of the chosen basis of the subspace
\(\Ran X\): if \(M\) is invertible, then \(P[XM]=P[X]\).

\section{Fixed-Charge Gaussian States as Occupied Frames}

\begin{definition}[Occupied frame]
A charge-\(q\) Slater state is represented by an \(N\times q\) frame
\(U=(u_1,\ldots,u_q)\) through
\begin{equation}
  \ket{U}
  =
  \ahat^\dagger(u_1)\cdots \ahat^\dagger(u_q)\ket{\vac}.
  \label{eq:slater-frame}
\end{equation}
The occupied one-particle subspace is \(\Ran U\), and its projector is
\begin{equation}
  C=P[U]=U(U^\dagger U)^{-1}U^\dagger .
  \label{eq:occupied-projector}
\end{equation}
\end{definition}

If \(U\mapsto UM\) with \(M\in GL(q,\mathbb C)\), the subspace and projector
do not change, while
\begin{equation}
  \ket{UM}=\det(M)^*\,\ket{U}.
\end{equation}
The physical normalized Gaussian state is therefore the occupied subspace,
whereas the frame stores a choice of basis and normalization.

The hole subspace and complex covariance are
\begin{equation}
  \cH_{\rm occ}=\Ran C,\qquad
  \cH_{\rm hole}=\Ker C=\Ran(\Id-C),
\end{equation}
\begin{equation}
  G=2C-\Id .
  \label{eq:complex-covariance}
\end{equation}
For a pure number-conserving Gaussian state,
\begin{equation}
  C^\dagger=C,\qquad C^2=C,\qquad G^\dagger=G,\qquad G^2=\Id .
\end{equation}

\subsection{State--Covariance Identities}

The fundamental Pauli constraints are
\begin{equation}
  \ahat^\dagger(Cx)\ket{U}=0,\qquad
  \ahat((\Id-C)x)\ket{U}=0 .
  \label{eq:frame-pauli}
\end{equation}
Creation in an already occupied direction vanishes; annihilation in an empty
direction vanishes.  In the local basis these become
\begin{align}
  \chat_i^\dagger\ket{G}
  &=
  -\sum_j G_{ij}\chat_j^\dagger\ket{G},
  \label{eq:creation-cov-identity}
  \\
  \chat_i\ket{G}
  &=
  \sum_j G_{ji}\chat_j\ket{G}.
  \label{eq:annihilation-cov-identity}
\end{align}
Equivalently,
\begin{equation}
  (\Id+G)\chat^\dagger\ket{G}=0,\qquad
  (\Id-G^T)\chat\ket{G}=0 .
\end{equation}
Thus the familiar covariance identities are not separate assumptions: they
are the occupied-frame Pauli constraints in component form.

\section{Chronological Frame Propagation}

Let \(C=P[U]\) be the current occupied projector.  The string is processed in
chronological order: the rightmost operator in a product acts first.

\subsection{Background State Updates}

\paragraph*{Gain.}
For a gain \(\ahat^\dagger(g)\), only the empty component can be added:
\begin{equation}
  h=(\Id-C)g.
\end{equation}
If \(h=0\), the state vanishes.  Otherwise
\begin{equation}
  C\longmapsto C+P[h],
  \label{eq:gain-background}
\end{equation}
or, at the frame level, append \(h\) and re-orthonormalize if desired.

\paragraph*{Loss.}
For a loss \(\ahat(\ell)\), only the occupied component can be removed:
\begin{equation}
  u=C\ell.
\end{equation}
If \(u=0\), the state vanishes.  Otherwise
\begin{equation}
  C\longmapsto C-P[u]
  =
  C-\frac{C\ell\,\ell^\dagger C}{\ell^\dagger C\ell}.
  \label{eq:loss-background}
\end{equation}
Geometrically, the new occupied subspace is
\begin{equation}
  \Ran C\cap \ell^\perp .
\end{equation}

\paragraph*{Even Gaussian.}
Let \(K\) be an invertible number-conserving Gaussian operator with
one-particle occupied-frame map \(R\):
\begin{equation}
  K\,\ahat^\dagger(v)\,K^{-1}=\ahat^\dagger(Rv).
  \label{eq:even-R}
\end{equation}
Then
\begin{equation}
  U\longmapsto RU,\qquad
  C\longmapsto P[RU].
  \label{eq:even-background}
\end{equation}
The associated annihilation-frame map is \(R^{-\dagger}\).

\subsection{Particle and Hole Tangents}

A particle tangent over the background \(C\) is an extra vector
\(r\in\Ran(\Id-C)\), representing
\(\ahat^\dagger(r)\ket{U}\).  After every update it is convenient to choose
the representative orthogonal to the updated occupied subspace.

For an even \(K\), with updated background \(C'\),
\begin{equation}
  r\longmapsto (\Id-C')Rr.
  \label{eq:particle-even}
\end{equation}
For a gain, after adding \(h=(\Id-C)g\),
\begin{equation}
  r\longmapsto (\Id-P[h])r.
  \label{eq:particle-gain}
\end{equation}
For a loss, with \(u=C\ell\),
\begin{equation}
  r\longmapsto r-u\,\frac{\ell^\dagger r}{\ell^\dagger u}.
  \label{eq:particle-loss}
\end{equation}
Equation~\eqref{eq:particle-loss} is an oblique projection: it subtracts the
removed occupied direction \(u\) so that the resulting vector is orthogonal
to the loss wavefunction \(\ell\).

The hole algorithm is the mirror construction on the empty projector
\(E=\Id-C\).  A hole tangent \(z\) represents \(\ahat(z)\ket{U}\) and is
kept in the occupied complement of the empty frame.  Under an even Gaussian
the hole wavefunction transforms by \(R^{-\dagger}\).  A loss adds an empty
cap, while a gain removes an empty direction obliquely.  In words,
\begin{equation}
  \text{particle: gains add caps, losses obliquely project,}
\end{equation}
\begin{equation}
  \text{hole: losses add caps, gains obliquely project.}
\end{equation}

\section{Choi Occupied Frames}

\subsection{Doubled Space and Reference Frame}

The Choi one-particle space is
\begin{equation}
  \cH_\Sigma=\cH_L\oplus\cH_R .
\end{equation}
A doubled vector is written as
\begin{equation}
  \binom{x_L}{x_R}.
\end{equation}
The number-conserving reference state is represented, up to an irrelevant
normalization, by the occupied frame
\begin{equation}
  F_{\REF}=
  \binom{\Id}{\Id}.
  \label{eq:ref-frame}
\end{equation}
The normalized projector would be
\begin{equation}
  C_{\REF}=\frac12
  \begin{pmatrix}
    \Id & \Id\\
    \Id & \Id
  \end{pmatrix},
\end{equation}
but the unnormalized frame in Eq.~\eqref{eq:ref-frame} gives the same
subspace.

For a right-leg string \(S_R\), define
\begin{equation}
  \ket{\Sigma_S}=S_R\ket{\REF}.
\end{equation}
If the final occupied frame is
\begin{equation}
  F=\binom{A}{B},
\end{equation}
then the Choi one-body projector and covariance are
\begin{equation}
  C_\Sigma=F(F^\dagger F)^{-1}F^\dagger,\qquad
  \Sigma=2C_\Sigma-\Id_{2N}.
  \label{eq:choi-projector}
\end{equation}

\subsection{Right-Leg Updates}

A right-leg gain \(\ahat_R^\dagger(g)\) is embedded as
\begin{equation}
  \tilde g=\binom{0}{g}.
\end{equation}
It appends the empty component
\begin{equation}
  h=(\Id-C_\Sigma)\tilde g
\end{equation}
provided \(h\neq 0\).

A right-leg loss \(\ahat_R(\ell)\) is embedded as
\begin{equation}
  \tilde\ell=\binom{0}{\ell}.
\end{equation}
At the subspace level, the loss keeps
\begin{equation}
  \Ran F\longmapsto \Ran F\cap \tilde\ell^\perp .
  \label{eq:choi-loss-intersection}
\end{equation}
Equivalently, if \(m=\tilde\ell^\dagger F\), choose any matrix \(K\) whose
columns span \(\Ker m\) and set \(F\mapsto FK\).  The choice-free projector
update is
\begin{equation}
  C_\Sigma\longmapsto
  C_\Sigma
  -
  \frac{C_\Sigma\tilde\ell\,\tilde\ell^\dagger C_\Sigma}
       {\tilde\ell^\dagger C_\Sigma\tilde\ell}.
  \label{eq:choi-loss-projector}
\end{equation}

A right-leg even Gaussian \(K_R\) with occupied-frame map \(R\) acts by
\begin{equation}
  F\longmapsto
  \begin{pmatrix}
    \Id & 0\\
    0 & R
  \end{pmatrix}F.
  \label{eq:choi-even}
\end{equation}

\section{Transfer Matrices from Choi Frames}

\subsection{Covariance Relation Pencils}

Block the Choi covariance as
\begin{equation}
  \Sigma=
  \begin{pmatrix}
    \Sigma_{LL} & \Sigma_{LR}\\
    \Sigma_{RL} & \Sigma_{RR}
  \end{pmatrix},
  \qquad
  C_\Sigma=
  \begin{pmatrix}
    C_{LL} & C_{LR}\\
    C_{RL} & C_{RR}
  \end{pmatrix}.
\end{equation}
The doubled state identities give the particle and hole relation pencils
\begin{align}
  A_p\chat_L^\dagger\ket{\Sigma}
  &=
  -B_p\chat_R^\dagger\ket{\Sigma},
  &
  A_p&=\Id+\Sigma_{LL},
  &
  B_p&=\Sigma_{LR},
  \label{eq:particle-pencil}
  \\
  A_h\chat_L\ket{\Sigma}
  &=
  B_h\chat_R\ket{\Sigma},
  &
  A_h&=\Id-\Sigma_{LL}^T,
  &
  B_h&=\Sigma_{RL}^T.
  \label{eq:hole-pencil}
\end{align}
When the relevant left block is invertible,
\begin{align}
  T_p
  &=
  (\Id+\Sigma_{LL})^{-1}\Sigma_{LR}
  =
  C_{LL}^{-1}C_{LR},
  \label{eq:Tp-cov}
  \\
  T_h
  &=
  (\Id-\Sigma_{LL}^T)^{-1}\Sigma_{RL}^T
  =
  (\Id-C_{LL})^{-T}C_{RL}^T .
  \label{eq:Th-cov}
\end{align}
If the inverse is singular, the pencil \((A_p,B_p)\) or \((A_h,B_h)\)
remains the exact Choi-covariance data.  A pseudoinverse is not a substitute
for the missing source or cap sector.

\subsection{Occupied Graph Coordinates and Particle Transfer}

Let
\begin{equation}
  F=\binom{A}{B}
\end{equation}
be a final occupied frame.  The right-only occupied cap is
\begin{equation}
  \cY=\{y\in\cH_R:\binom{0}{y}\in\Ran F\}.
  \label{eq:Y-definition}
\end{equation}
Algebraically, if \(Y\) is any frame for this subspace, then
\begin{equation}
  \Ran Y=B\,\Ker A .
  \label{eq:Y-BkerA}
\end{equation}
When \(A\) has full row rank, choose any right inverse \(J\) satisfying
\begin{equation}
  AJ=\Id .
\end{equation}
Then
\begin{equation}
  FJ=\binom{\Id}{BJ},
\end{equation}
so \(R_0=BJ\) is a graph representative.  It is not unique: changing \(J\)
changes \(R_0\) by a map into \(\cY\).  The covariance chart selects the
canonical representative
\begin{equation}
  R=Q[Y]R_0.
  \label{eq:canonical-R}
\end{equation}
The particle transfer matrix in the Choi convention is
\begin{equation}
  T_p=R^\dagger.
  \label{eq:Tp-frame}
\end{equation}
Equivalently, \(T_p\) is the unique representative of the particle relation
with
\begin{equation}
  T_pY=0.
  \label{eq:Tp-kills-Y}
\end{equation}
The filled cap \(Y\) is not a choice; it is the right-only occupied part of
the Choi subspace.  Only the extension of the transfer matrix on that killed
subspace is a coordinate choice.

\subsection{Empty Graph Coordinates and Hole Transfer}

Let \(E=\binom{A_h}{B_h}\) be a frame for the empty Choi subspace
\(\Ran(\Id-C_\Sigma)\).  The right-only empty cap is
\begin{equation}
  \cZ=\{z\in\cH_R:\binom{0}{z}\in\Ran E\},
  \qquad
  \Ran Z=B_h\Ker A_h .
  \label{eq:Z-definition}
\end{equation}
If \(A_h\) has full row rank, choose \(J_h\) with \(A_hJ_h=\Id\), set
\begin{equation}
  H_0=B_hJ_h,\qquad H=Q[Z]H_0.
\end{equation}
For an empty graph \(E\sim \binom{\Id}{H}\), Eq.~\eqref{eq:Th-cov} gives
\begin{equation}
  T_h=-H^T .
  \label{eq:Th-frame}
\end{equation}
The canonical hole representative kills the conjugated empty cap,
\begin{equation}
  T_h\overline{Z}=0.
\end{equation}
The conjugation is a consequence of the hole convention in
Eq.~\eqref{eq:hole-pencil}.

\section{Regulated Endcaps}

Suppose the finite particle chart has the form
\begin{equation}
  T_p=N_pQ[Y],
\end{equation}
where \(Y\) is the filled right-output cap.  A simple endcap regulator is
\begin{equation}
  Q[Y]\longmapsto Q[Y]+\eps P[Y],
\end{equation}
so
\begin{equation}
  T_{p,\eps}=N_p\bigl(Q[Y]+\eps P[Y]\bigr).
  \label{eq:particle-regulator}
\end{equation}
The reciprocal hole representative is
\begin{equation}
  T_{h,\eps}=T_{p,\eps}^{-T}.
  \label{eq:hole-reciprocal}
\end{equation}
Thus an exact zero in the particle chart appears as a reciprocal pole in the
hole representative.

For a finite hole chart with empty cap \(Z\),
\begin{equation}
  T_h=N_hQ[\overline Z],
\end{equation}
the mirror regulator is
\begin{equation}
  T_{h,\eps}=N_h\bigl(Q[\overline Z]+\eps P[\overline Z]\bigr),
  \qquad
  T_{p,\eps}=T_{h,\eps}^{-T}.
  \label{eq:hole-regulator}
\end{equation}
The cap subspace is physical Choi-covariance data.  The finite
\(\eps\)-extension is a regulated square-matrix representative.

\section{Worked Examples}

\subsection{\texorpdfstring{One Gain Between Evens: \(K_{2,R}\ahat_R^\dagger(\chi)K_{1,R}\ket{\REF}\)}{One Gain Between Evens}}

Let
\begin{equation}
  K_i\ahat^\dagger(v)K_i^{-1}=\ahat^\dagger(R_i v),
  \qquad
  T_i=R_i^\dagger .
\end{equation}
Starting from \(F_{\REF}\), after \(K_{1,R}\) one has
\begin{equation}
  F_1=\binom{\Id}{R_1}.
\end{equation}
The right gain appends \(\binom{0}{\chi}\), modulo the already occupied graph:
\begin{equation}
  F_2\sim
  \begin{pmatrix}
    \Id & 0\\
    R_1 & \chi
  \end{pmatrix}.
\end{equation}
After \(K_{2,R}\),
\begin{equation}
  F_f\sim
  \begin{pmatrix}
    \Id & 0\\
    R_2R_1 & R_2\chi
  \end{pmatrix}.
\end{equation}
Define
\begin{equation}
  y=R_2\chi=T_2^\dagger\chi,\qquad Q_y=Q[y].
\end{equation}
The canonical occupied graph is
\begin{equation}
  R_{\rm can}=Q_yR_2R_1.
\end{equation}
Hence
\begin{equation}
  T_p=R_{\rm can}^\dagger=T_1T_2Q_y,
\end{equation}
and the regulated particle and reciprocal hole representatives are
\begin{equation}
  T_{p,\eps}=T_1T_2(Q_y+\eps P[y]),
  \qquad
  T_{h,\eps}=T_{p,\eps}^{-T}.
  \label{eq:single-gain-regulated}
\end{equation}
There is a local Pauli projection at the gain insertion point,
\(Q[\chi]\), in the chronological particle tangent.  It is redundant in the
canonical final representative because
\begin{equation}
  Q_yR_2P[\chi]=0.
\end{equation}

\subsection{\texorpdfstring{Gain Followed by Loss: \(K_{3,R}\ahat_R(\ell)K_{2,R}\ahat_R^\dagger(\chi)K_{1,R}\ket{\REF}\)}{Gain Followed by Loss}}

After \(K_1\), the gain, and \(K_2\), write
\begin{equation}
  F_-=
  \begin{pmatrix}
    \Id & 0\\
    A & y
  \end{pmatrix},
  \qquad
  A=R_2R_1,\qquad y=R_2\chi.
  \label{eq:pre-loss-frame}
\end{equation}
A vector in \(\Ran F_-\) has the form
\begin{equation}
  \binom{x}{Ax+\alpha y}.
\end{equation}
The right loss by \(\ell\) imposes the condition
\begin{equation}
  \binom{0}{\ell}^\dagger
  \binom{x}{Ax+\alpha y}
  =
  \ell^\dagger(Ax+\alpha y)=0.
  \label{eq:loss-alpha-condition}
\end{equation}
If
\begin{equation}
  s=\ell^\dagger y\neq 0,
\end{equation}
then
\begin{equation}
  \alpha=-s^{-1}\ell^\dagger Ax.
  \label{eq:alpha-solution}
\end{equation}
Therefore
\begin{equation}
  Ax+\alpha y
  =
  \left[\Id-y(\ell^\dagger y)^{-1}\ell^\dagger\right]Ax.
\end{equation}
The bracket is the oblique projector
\begin{equation}
  O_{y|\ell}=\Id-y(\ell^\dagger y)^{-1}\ell^\dagger,
  \qquad
  O_{y|\ell}y=0,\qquad
  \ell^\dagger O_{y|\ell}=0.
\end{equation}
After \(K_{3,R}\),
\begin{equation}
  R_f=R_3O_{y|\ell}R_2R_1.
\end{equation}
Thus
\begin{equation}
  T_p
  =
  T_1T_2
  \left[\Id-\ell(y^\dagger\ell)^{-1}y^\dagger\right]T_3.
  \label{eq:gain-loss-Tp}
\end{equation}
The regulated form is
\begin{equation}
  T_{p,\eps}
  =
  T_1T_2
  \left[
    \Id-(1-\eps)\ell(y^\dagger\ell)^{-1}y^\dagger
  \right]T_3,
  \qquad
  T_{h,\eps}=T_{p,\eps}^{-T}.
  \label{eq:gain-loss-regulated}
\end{equation}
No canonical assumption was used to determine \(\alpha\): it is fixed by the
loss intersection condition \(\Ran F_-\cap (0\oplus\ell)^\perp\).

\subsection{\texorpdfstring{Two Consecutive Gains: \(K_{2,R}\ahat_R^\dagger(\chi_1)\ahat_R^\dagger(\chi_2)K_{1,R}\ket{\REF}\)}{Two Consecutive Gains}}

The rightmost gain \(\chi_2^\dagger\) acts first.  After \(K_{2,R}\), the
occupied frame can be written
\begin{equation}
  F_f\sim
  \begin{pmatrix}
    \Id & 0 & 0\\
    R_2R_1 & R_2\chi_2 & R_2\chi_1
  \end{pmatrix}.
\end{equation}
Let
\begin{equation}
  Y=R_2(\chi_2,\chi_1)=T_2^\dagger(\chi_2,\chi_1).
\end{equation}
Assuming the two gain directions remain independent after projection, the
finite particle representative is
\begin{equation}
  T_p=T_1T_2Q[Y],
\end{equation}
and the regulated pair is
\begin{equation}
  T_{p,\eps}=T_1T_2\bigl(Q[Y]+\eps P[Y]\bigr),
  \qquad
  T_{h,\eps}=T_{p,\eps}^{-T}.
  \label{eq:two-gains-regulated}
\end{equation}

\subsection{\texorpdfstring{Two Consecutive Losses: \(K_{2,R}\ahat_R(\chi_2)\ahat_R(\chi_1)K_{1,R}\ket{\REF}\)}{Two Consecutive Losses}}

For losses the empty frame is the natural object.  After \(K_{1,R}\), an
empty frame orthogonal to \(\binom{\Id}{R_1}\) is
\begin{equation}
  E_1=
  \binom{\Id}{-R_1^{-\dagger}}.
\end{equation}
The two losses add right-only empty directions.  After \(K_{2,R}\), empty
wavefunctions are propagated by \(R_2^{-\dagger}\), giving
\begin{equation}
  E_f\sim
  \begin{pmatrix}
    \Id & 0 & 0\\
    -R_2^{-\dagger}R_1^{-\dagger}
      & R_2^{-\dagger}\chi_1
      & R_2^{-\dagger}\chi_2
  \end{pmatrix}.
\end{equation}
Define the empty cap frame
\begin{equation}
  Z=R_2^{-\dagger}(\chi_1,\chi_2).
\end{equation}
The finite hole representative is
\begin{equation}
  T_h=T_1^{-T}T_2^{-T}Q[\overline Z],
\end{equation}
with regulated pair
\begin{equation}
  T_{h,\eps}
  =
  T_1^{-T}T_2^{-T}
  \bigl(Q[\overline Z]+\eps P[\overline Z]\bigr),
  \qquad
  T_{p,\eps}=T_{h,\eps}^{-T}.
  \label{eq:two-losses-regulated}
\end{equation}

\subsection{A Two-Mode Sanity Check}

Set all evens to the identity, choose \(N=2\), let
\begin{equation}
  \chi=e_1,\qquad \ell=e_1+\alpha e_2 .
\end{equation}
For the gain--loss word \(\ahat(\ell)\ahat^\dagger(\chi)\), one has
\begin{equation}
  y=e_1,\qquad y^\dagger\ell=1,
\end{equation}
and
\begin{equation}
  \ell(y^\dagger\ell)^{-1}y^\dagger
  =
  \begin{pmatrix}
    1 & 0\\
    \alpha & 0
  \end{pmatrix}.
\end{equation}
Equation~\eqref{eq:gain-loss-regulated} gives
\begin{equation}
  T_{p,\eps}
  =
  \Id-(1-\eps)\ell y^\dagger
  =
  \begin{pmatrix}
    \eps & 0\\
    -(1-\eps)\alpha & 1
  \end{pmatrix}.
\end{equation}
The finite limit
\begin{equation}
  T_p=
  \begin{pmatrix}
    0 & 0\\
    -\alpha & 1
  \end{pmatrix}
\end{equation}
has a particle zero in the filled direction, while
\begin{equation}
  T_{h,\eps}=T_{p,\eps}^{-T}
\end{equation}
has the reciprocal pole.

\section{Summary}

The occupied-frame construction separates three pieces of data that are often
conflated in transfer-matrix notation.
First, the Choi covariance is the primary state-level object:
\(C_\Sigma=P[F]\).  Second, the finite particle or hole transfer matrix is a
graph coordinate of the occupied or empty Choi subspace.  Third, a non-neutral
string creates an exact right-output cap: a filled cap \(Y\) for net gain or
an empty cap \(Z\) for net loss.  The covariance chart returns the canonical
representative that kills this cap, while the regulated endcap
\(Q+\eps P\) supplies an invertible square representative whose opposite
particle--hole sector has a reciprocal pole.

\end{document}
